\DTMsavedate{mydate}{2022-06-30}
\maketitle{Electrical Engineering - Basic Circuit Analysis}{Electronic Engineering}{\DTMusedate{mydate}}

% ----------------------------------------------------- %
% COURSE INFO
\subsection*{Course Information}
\vspace{0ex}%
\noindent\parbox{0.5\textwidth}{%
    \noindent\begin{tabular}{@{}ll}
                 \textsf{Course:}          & Electrical Engineering - Basic Circuit Analysis \\
                 \textsf{Course Code:}         & 	TEE 1103    \\
                 \textsf{Credit Hours:}    & 	10
    \end{tabular}}
\vspace{2ex}\hrule\vspace{2ex}

% ----------------------------------------------------- %
% INSTRUCTOR INFO
\subsection*{Lecturer Information}
\noindent\parbox{0.5\textwidth}{%
    \noindent\begin{tabular}{@{}ll}
                 \textsf{Name:}         &    Svetlana A. Bebova           \\
% \textsf{Phone:}        &    \phone          \\
\textsf{Email:}        &    \href{mailto:swetlana.bebova@nust.ac.zw}{swetlana.bebova@nust.ac.zw} \\
\textsf{Qualifications:}        &  MSc in Radio Electronic Engineering (Technical University of Sofia)
    \end{tabular}}
\vspace{2ex}\hrule\vspace{2ex}


% ----------------------------------------------------- %
% COURSE SYNOPSIS
\subsection*{Course Synopsis}
The course introduces the student to general concepts of current, voltage, resistance and circuits (dc and ac).
It also covers elements of loop and nodal analysis, basic networks and theorems, Delta-Wye conversions and network theorems.
It delves into Capacitor and inductor circuits, Transient analysis, DC and AC power.
It also delves into forced response and sinusoidal steady-state response.
\vspace{2ex} \hrule\vspace{2ex}
% ----------------------------------------------------- %
% COURSE TOPICS
\subsection*{Topics}
\begin{outline}
    \1 Introduction and Review of Electricity
    \1 Basics of Circuit Analysis (DC)
    \1 Methods of Circuit Analysis (DC)
    \1 Network Theorem (DC)
    \1 Capacitance
    \1 Magnetic Circuits
    \1 Inductance
    \1 Alternating Current
    \1 Analysis of AC Circuits
    \1 Application of Network Theorems to AC Circuits
    \1 AC Power
\end{outline}
\vspace{2ex} \hrule\vspace{2ex}

% Students Assenssment
\subsection*{Assessment of Students}
\begin{outline}
    \1 Students are assessed through written assignments and written in-class tests that shall form part of continuous assessment mark.
    \1 There shall be a final examination at the end of the semester.
\end{outline}
\vspace{2ex}\hrule\vspace{2ex}

% ----------------------------------------------------- %
% Grade Evaluation
\subsection*{Course Evaluation and Grading Policies}
Grades are based on:
Continuous assessment (\qty{25}{\percent}),
Final Exam (\qty{75}{\percent})
\vspace{2ex}\hrule\vspace{2ex}

% Policies and Other information
\subsection*{Policies and Other Information}
% Conduct
\subsection*{Key Text}
\begin{itemize}
    \item Introductory Circuit Analysis, 11th Edition, 2007 by R. L. Boylestad
    \item Introductory Electronic Devices and Circuits: Electron Flow Version, 7th Edition, 2006 by R. Paynter
    \item Electronic Measurement and Instrumentation, 1996 by K. B. Klaassen, S. Gee
\end{itemize}
\vspace{2ex}\hrule
\vspace{2ex}\hrule\vspace{2ex}